\subsection{交错多重线性形式的乘积}

设 A 是一个交换环，M 是一个 A-模。对每个整数 \(r \geq 0\)，用 \(\operatorname{Alt}^{r}(\mathbf{M})\) 表示 M 上交错 r-线性形式所成的 A-模；它可等同于 A-模 \(\Lambda^{r}(\mathbf{M})\) 的对偶（《代数》，第三章，§7，第 4 节，命题 7）。设 \(u \in \operatorname{Alt}^{s}(\mathbf{M})\) 与 \(v \in \operatorname{Alt}^{r}(\mathbf{M})\)；回忆（《代数》，第三章，§11，第 2 节，例 3）u 与 v 的交错积是由下式定义的元素 \(u \wedge v \in \operatorname{Alt}^{s+r}(\mathbf{M})\)：

\[
(u \wedge v) (x _ {1}, \dots , x _ {s + r}) = \sum_ {\sigma \in \mathfrak {S} _ {s, r}} \varepsilon_ {\sigma} u (x _ {\sigma (1)}, \dots , x _ {\sigma (s)}) v (x _ {\sigma (s + 1)}, \dots , x _ {\sigma (s + r)}),
\]

其中 \(S_{s,r}\) 是对称群 \(S_{s+r}\) 中由在 \((1,s)\) 与 \((s+1,s+r)\) 上的限制都是递增的置换所组成的子集。

现在设

\[
0 \longrightarrow \mathrm{M} ^ {\prime} \stackrel {i} {\longrightarrow} \mathrm{M} \stackrel {p} {\longrightarrow} \mathrm{M} ^ {\prime \prime} \longrightarrow 0
\]

是自由 A-模的一个正合序列，其秩分别为 \(r, r + s\) 与 \(s\)。

引理 1. 存在一个从 \(\mathrm{Alt}^s (\mathbf{M}'')\times \mathrm{Alt}^r (\mathbf{M}')\) 到 \(\mathrm{Alt}^{s + r}(\mathbf{M})\) 的 A-双线性映射，记作 \((u,v)\mapsto u\cap v\)，它由下列两个性质之一刻画：

\begin{enumerate}
\def\labelenumi{\alph{enumi})}
\item
  用 \(u_{1} \in \operatorname{Alt}^{s}(\operatorname{M})\) 表示形式 \((x_{1}, \ldots, x_{s}) \mapsto u(p(x_{1}), \ldots, p(x_{s}))\)，并设 \(v_{1} \in \operatorname{Alt}^{r}(\operatorname{M})\) 是一个满足 \(v_{1}(i(x_{1}^{\prime}), \ldots, i(x_{r}^{\prime})) = v(x_{1}^{\prime}, \ldots, x_{r}^{\prime})\)（对 \(x_{1}^{\prime}, \ldots, x_{r}^{\prime}\)，\(M^{\prime}\) 中的）的形式；则 \(u \cap v = u_{1} \wedge v_{1}\)。
\item
  对所有 \(x_{1},\ldots,x_{s}\)（M 中）与 \(x_{1}^{\prime},\ldots,x_{r}^{\prime}\)（\(M^{\prime}\) 中），
\end{enumerate}

\[
(u \cap v) (x _ {1}, \dots , x _ {s}, i (x _ {1} ^ {\prime}), \dots , i (x _ {r} ^ {\prime})) = u (p (x _ {1}), \dots , p (x _ {s})) v (x _ {1} ^ {\prime}, \dots , x _ {r} ^ {\prime}).\tag{1}
\]

映射 \(\varphi : \mathrm{Alt}^s (\mathrm{M}'')\otimes_{\mathrm{A}}\mathrm{Alt}^r (\mathrm{M}')\to \mathrm{Alt}^{s + r}(\mathrm{M})\)（满足 \(\varphi (u\otimes v) = u\cap v\)）是秩 1 自由 A-模之间的一个同构。

形式 \(v_{1}\)（满足条件 \(a\))）的存在性来自下述事实：\(\Lambda^r (i)\) 诱导出从 \(\Lambda^r (\mathrm{M}')\) 到 \(\Lambda^r (\mathrm{M})\) 的一个直因子子模的同构（《代数》，第三章，§7，第 2 节）。设 \(v_{1}\) 是这样一个形式；令 \(u\cap v = u_{1}\wedge v_{1}\)。这时公式 (1) 得到满足，因为若令 \(i(x_k') = x_{s + k}\)（对 \(1\leq k\leq r\)），则元素 \(\sigma\)（属于 \(\mathfrak{S}_{s,r}\)）中使 \(p(x_{\sigma (i)})\neq 0\)（对 \(1\leq i\leq s\)）的只有恒等置换。另一方面，公式 (1) 唯一确定 \(u\cap v\)：实际上，设 \((e_1',\ldots ,e_r')\) 是 \(\mathrm{M}'\) 的一个基，\((f_1'',\dots,f_s'')\) 是 \(\mathrm{M}''\) 的一个基，而 \(f_{1},\ldots ,f_{s}\) 是 \(\mathrm{M}\) 中满足 \(p(f_i) = f_i''\)（对 \(1\leq i\leq s\)）的元素。则 \((f_{1},\ldots ,f_{s},i(e_1'),\ldots ,i(e_r'))\) 是 \(\mathrm{M}\) 的一个基（《代数》，第二章，§1，第 11 节，命题 21），而公式 (1) 可写成

\[
(u \cap v) (f _ {1}, \dots , f _ {s}, i (e _ {1} ^ {\prime}), \dots , i (e _ {r} ^ {\prime})) = u (f _ {1} ^ {\prime \prime}, \dots , f _ {s} ^ {\prime \prime}) v (e _ {1} ^ {\prime}, \dots , e _ {r} ^ {\prime});\tag{2}
\]

但 \(\mathrm{Alt}^{s + r}(\mathbf{M})\) 的一个元素由它在一个基上的值唯一确定。

由前述可知，条件 a) 与 b) 中的每一个都唯一确定乘积 \(u \cap v\)；显然这个乘积是双线性的。最后，引理的最后一个断言由公式 (2) 得出。
% semantic-fragments: legacy-input-seam
\relax{}
